\documentclass[11pt,aspectratio=169]{beamer}

% ---------- Theme ----------
\usetheme{Madrid}
\usecolortheme{seahorse}
\setbeamertemplate{navigation symbols}{}
\setbeamertemplate{footline}[frame number]
\setbeamertemplate{caption}[numbered]

\usepackage[utf8]{inputenc}
\usepackage[T1]{fontenc}
\usepackage{booktabs}
\usepackage{array}
\usepackage{amsmath}
\usepackage{tikz}
\usetikzlibrary{arrows.meta, positioning}

\newcolumntype{L}[1]{>{\raggedright\arraybackslash}p{#1}}

\title[SLURRY\_FREE\_ACID Progress]{SLURRY\_FREE\_ACID Modeling and Process Intelligence}
\subtitle{Project Progress and Model Evaluation}
\author{Ikram Benfellah}
\institute{OCP Group — Internship Progress Update}
\date{\today}

\begin{document}

% =====================================================================
\begin{frame}
\titlepage
\end{frame}
\note{
Quick framing before I start: this is a working update, not a defense. I want to walk through what I built, what I found, what worked, what didn't, and what I think we should do next.
}


% =====================================================================
\begin{frame}{Dataset and Evaluation Design}
\begin{columns}[T]
\begin{column}{0.5\textwidth}
\textbf{Dataset}
\begin{itemize}
    \small
    \item 44{,}640 chronological observations
    \item $\approx$31 days, January 2026
    \item One observation per minute
    \item 21 original process measurements
    \item Target: \texttt{SLURRY\_FREE\_ACID}
\end{itemize}
\end{column}
\begin{column}{0.5\textwidth}
\textbf{Chronological split}
\begin{itemize}
    \small
    \item $\approx$70\% train / 15\% validation / 15\% test
    \item No random shuffling
    \item All preprocessing (imputation, scaling, lag selection, filtering, clipping) fitted on training data only
\end{itemize}
\end{column}
\end{columns}
\vspace{0.3cm}
\begin{alertblock}{Why chronological, not random}
\small The 44,640 rows are \textbf{not} 44,640 independent experiments — they are one continuous month with strongly dependent adjacent minutes.
\end{alertblock}
\begin{block}{Feature preparation (brief)}
\small I built four feature sets as the experimental foundation: A — Raw (21 features), B — Temporal (162 retained), C — Process-aware (58 retained), D — Full (178 retained). I won't detail the engineering here — it's covered in the notebooks.
\end{block}
\end{frame}
\note{
Two things to stress here: chronological splitting was a deliberate choice to avoid leaking near-identical minutes across train/test, and the one-month limitation is structural, not something I can fix by adding more rows from the same month. I'll only spend a few seconds on the four feature sets — happy to go deeper offline if useful.
}

% =====================================================================
\begin{frame}{Three Operational Prediction Cases}
This distinction matters for how any of these models could actually be used.
\vspace{0.2cm}
\begin{center}
\small
\begin{tabular}{L{2.4cm} L{3.1cm} L{2.6cm} L{2.4cm} L{2.6cm}}
\toprule
\textbf{Case} & \textbf{Required input} & \textbf{Predicts} & \textbf{Purpose} & \textbf{Main limitation} \\
\midrule
1. Direct forecast & Process measurements up to $t$ & $y(t+h)$ & Early warning, $h$=1/5/10/15 min & Must beat persistence to be useful \\
\addlinespace
2. Delta forecast & Process measurements $+$ current/past $y$ & $\Delta y(t,h)$, reconstructed as $y(t)+\Delta y$ & Short-horizon early warning & Needs $y(t)$ at inference \\
\addlinespace
3. Virtual sensor & Process measurements only & $y(t)$ (current, estimated) & Soft sensor when no lab/analyzer value exists & Estimates \emph{now}, not the future \\
\bottomrule
\end{tabular}
\end{center}
\end{frame}
\note{
This slide is the key to avoid confusion later: Case 1 and 2 forecast the future and need different inputs; Case 3 estimates the current value and is a completely different use case — it's a virtual sensor, not an early-warning tool. I'll refer back to this table for the rest of the deck.
}

% =====================================================================
\begin{frame}{Direct Forecasting and the Persistence Result}
\textbf{Notebook 03 — all 16 combinations (4 horizons $\times$ 4 feature sets)}
\begin{itemize}
    \item Baselines: training-target mean, and persistence ($\hat{y}(t+h) = y(t)$) when $y(t)$ was available.
    \item Models tested: Ridge, Extra Trees, Histogram Gradient Boosting.
\end{itemize}
\begin{block}{Result}
The direct future-level models learned the general operating level of SLURRY\_FREE\_ACID, but did \textbf{not consistently outperform persistence}.
\end{block}
\begin{exampleblock}{Why}
SLURRY\_FREE\_ACID changes relatively slowly over much of the month, so the current value $y(t)$ is already a very strong predictor at short horizons.
\end{exampleblock}
\textbf{Consequence:} this motivated predicting the \emph{change} instead of the full future level — the approach in Notebook 04.
\end{frame}
\note{
This was a useful negative result, not a failure — it told me the problem wasn't "learn the future level from scratch," it was "learn how much it moves." That reframing is what Notebook 04 is built on.
}

% =====================================================================
\begin{frame}{Target-Anchored Delta Forecasting (Notebook 04)}

\begin{columns}[T]

\begin{column}{0.55\textwidth}

\textbf{Reformulation}

\[
\Delta y(t,h) = y(t+h) - y(t)
\]

\begin{itemize}
    \small
    \item Model learns only the movement.
    \item Best candidate: \textbf{Ridge}, hybrid/full feature configuration
    (process measurements + target history).
\end{itemize}

\end{column}

\begin{column}{0.45\textwidth}

\begin{exampleblock}{Example}
\small

current $y(t) = 11.20$ \\

future $y(t+h) = 11.35$ \\

model only needs to learn: $+0.15$

\end{exampleblock}

\end{column}

\end{columns}

\vspace{0.2cm}

\begin{block}{Result}
This model beat persistence at \textbf{all four horizons}, with an improvement
of roughly \textbf{19.8\%--34.9\%} across horizons.
\end{block}

\begin{alertblock}{Limitation}
\small
Requires $y(t)$ at inference --- cannot be used on a period with no
\texttt{SLURRY\_FREE\_ACID} measurement available at all.
\end{alertblock}

\end{frame}
\note{
The intuition worth repeating out loud: most of the future value is already contained in the current value, so asking the model to predict the small movement is a much easier and more accurate task than asking it to predict the whole level again. This is the strongest result so far, but it only works when we already have a current reading.
}

% =====================================================================
\begin{frame}{Target-Free Virtual Sensor (Notebook 05)}
\begin{columns}[T]
\begin{column}{0.55\textwidth}
\textbf{Contract}
\begin{itemize}
    \small
    \item Training: $X$ = process features, $y$ = current SLURRY\_FREE\_ACID (target excluded from inputs).
    \item Inference: a new CSV with process measurements only can receive current free-acid estimates.
    \item Models evaluated: Ridge, PLS, KNN, Extra Trees, HGB, XGBoost, LightGBM, and weighted blends.
    \item Best: \textbf{50/50 Ridge + LightGBM blend} — Ridge for stable linear structure, LightGBM for nonlinear thresholds/interactions.
\end{itemize}
\end{column}
\begin{column}{0.45\textwidth}
\begin{block}{Retrospective test performance}
\small
\begin{tabular}{ll}
MAE & 0.2490 \\
RMSE & 0.3014 \\
R\textsuperscript{2} & 0.5583 \\
\end{tabular}
\vspace{0.15cm}
\small $R^2 \approx 0.558$ means the model explained about 55.8\% of target variance relative to a constant-mean reference — \emph{not} ``55.8\% accuracy.''
\end{block}
\end{column}
\end{columns}
\end{frame}
\note{
I want to be precise about R-squared here since it's easy to misread: it's variance explained relative to just predicting the mean, not an accuracy percentage. This is the baseline target-free model that everything in the next few slides tries to improve on.
}

% =====================================================================
\begin{frame}{Expanding Validation and the Notebook 06 Upgrade}
\textbf{Four expanding chronological validation blocks}, each $\approx$1,657 rows, training set growing each time (30{,}940 $\to$ 35{,}911 rows).
\vspace{0.15cm}

Candidates tested: stronger/global Ridge, recency-weighted Ridge (half-life 3/7/14 days), clustered Ridge experts (3 or 5 KMeans regions), LightGBM variants, Ridge–LightGBM blends.

\vspace{0.1cm}
\begin{center}
\scriptsize
\begin{tabular}{lccc}
\toprule
\textbf{Candidate} & \textbf{Mean RMSE} & \textbf{RMSE SD} & \textbf{Worst-block RMSE} \\
\midrule
70\% Ridge / 30\% LightGBM & 0.3191 & 0.0762 & 0.4236 \\
50\% Ridge / 50\% LightGBM & 0.3201 & 0.0686 & 0.4124 \\
\textbf{5-cluster Ridge (selected)} & 0.3280 & 0.0396 & 0.3522 \\
\bottomrule
\end{tabular}
\end{center}

\begin{block}{Selection and result}
\small I selected \textbf{5-cluster Ridge} as the stability-first candidate: it has the lowest validation RMSE SD (0.0396) and best worst-block RMSE (0.3522). This accepts weaker mean RMSE than the 70/30 blend. After refitting on train+validation, retrospective test: MAE 0.2354, RMSE 0.2977 ($\approx$1.23\% below Notebook 05), R\textsuperscript{2} 0.5692. The clusters are statistical operating regions, not confirmed physical regimes.
\end{block}
\end{frame}
\note{
Worth flagging explicitly: this improvement wasn't a new model, it was better methodology — a sturdier expanding validation scheme and stability-aware selection. I think that distinction matters when we talk about where future gains will come from.
}

% =====================================================================
\begin{frame}{Regime-Aware Hybrid: Accuracy vs. Stability (Notebook 07)}
Notebook 06 selected 5-cluster Ridge for stability. Notebook 07 combines it with the best-mean \textbf{70/30 Ridge--LightGBM} accuracy anchor:
\[
\hat{y} = (1-w)\cdot \hat{y}_{\text{global}} + w \cdot \hat{y}_{\text{local}}
\]
\begin{center}
\scriptsize
\begin{tabular}{lccc}
\toprule
\textbf{Configuration} & \textbf{Mean RMSE} & \textbf{RMSE SD} & \textbf{Worst-block RMSE} \\
\midrule
Global only ($w{=}0$) & 0.3191 & 0.0762 & 0.4236 \\
\textbf{80\% local (selected)} & 0.3213 & 0.0413 & 0.3548 \\
Local only ($w{=}1$) & 0.3280 & 0.0396 & 0.3522 \\
\bottomrule
\end{tabular}
\end{center}
\begin{block}{Selected: 80\% local hybrid, vs. global}
\small Against the 70/30 accuracy anchor, the selected hybrid accepts $\approx$0.69\% higher mean RMSE while reducing worst-block RMSE by $\approx$16.25\% and RMSE SD by $\approx$45.78\%. Retrospective test: MAE 0.2324, RMSE 0.2903, R\textsuperscript{2} 0.5902; RMSE is $\approx$2.47\% below Notebook 06's clustered model. \emph{This remains retrospective evidence.}
\end{block}
\end{frame}
\note{
This is a robustness-oriented alternative, not a strictly better model — it trades a small amount of average accuracy for meaningfully more consistent worst-case behavior across blocks. I'd frame the choice between Notebook 06 and 07 as "which failure mode do we care about more," not "which one is objectively better."
}

% =====================================================================
\begin{frame}{JITL Experiment — and Why It Was Rejected (Notebook 08)}
\textbf{Idea:} instead of 5 fixed clusters, find query-specific historical neighborhoods and fit a small local Ridge model per query, blended with the global model.
\begin{itemize}
    \small
    \item Similarity spaces: PCA (10/20 comp.), PLS (5/10 comp.)
    \item Neighbor counts: 100 / 250 / 500; weighting: uniform / inverse distance
    \item Local contribution: 25\% / 50\% / 75\% / 100\%
    \item 97 candidate configurations, evaluated across the same 4 validation blocks
\end{itemize}
\begin{center}
\scriptsize
\begin{tabular}{lccc}
\toprule
& \textbf{Mean RMSE} & \textbf{RMSE SD} & \textbf{Worst-block RMSE} \\
\midrule
Best JITL (PLS-10, 500 nb., uniform, 25\% local) & 0.3254 & 0.0799 & 0.4353 \\
70/30 global accuracy anchor & 0.3191 & 0.0762 & 0.4236 \\
\bottomrule
\end{tabular}
\end{center}
\begin{alertblock}{Decision: rejected}
\small Best JITL was $\approx$1.98\% worse on mean RMSE and $\approx$2.75\% worse on worst-block RMSE than the 70/30 global anchor, with higher variability and cost. Validation therefore retained the global anchor. Likely limits: similarity in measured variables may not imply the same free-acid relationship, and one month provides limited operating diversity.
\end{alertblock}
\end{frame}
\note{
I think this negative result is worth presenting on equal footing with the positive ones — I tested a genuinely more complex idea, and the chronological evidence didn't support it, so I didn't keep it just because it was more sophisticated.
}

% =====================================================================
\begin{frame}{Final Modeling Decision}
\begin{center}
\scriptsize
\begin{tabular}{L{2.6cm} L{3.4cm} L{4.2cm} L{1.8cm}}
\toprule
\textbf{Candidate} & \textbf{Role} & \textbf{Retrospective metrics} & \textbf{Status} \\
\midrule
Ridge delta model & Target-anchored forecasting & Beat persistence, all horizons, $\approx$19.8\%–34.9\% better & Selected \\
\addlinespace
5-cluster Ridge & Stability-first target-free & MAE 0.2354, RMSE 0.2977, $R^2$ 0.5692 & Selected \\
\addlinespace
80\% local regime hybrid & Balanced target-free alternative & MAE 0.2324, RMSE 0.2903, $R^2$ 0.5902 & Selected (alt.) \\
\addlinespace
JITL & Query-specific local model & Worse mean \& worst-block RMSE, higher variability & Rejected \\
\addlinespace
Operational model & Deployed / control-grade & — & \textbf{Not selected} — needs external validation on a later labeled month \\
\bottomrule
\end{tabular}
\end{center}
\end{frame}
\note{
This is the slide I'd want to leave on the screen the longest if we discuss priorities — three candidates are in good shape for their respective use cases, one was correctly discarded, and none of them is ready to be called an operational model yet.
}

% =====================================================================
\begin{frame}{Dashboard and Row-Level Reliability Outputs}
\begin{columns}[T]
\begin{column}{0.55\textwidth}
\textbf{Capabilities}
\begin{itemize}
    \small
    \item Default January dataset or CSV upload
    \item Process-variable exploration, correlations, relationship explorer
    \item Candidate transition detection (top 5\% multivariate change score — analytical label, not confirmed plant events)
    \item Target-free current-value prediction \& target-anchored forecasts (1/5/10/15 min)
    \item Downloadable predictions, drift/familiarity info, report generation
    \item Loads saved artifacts only — \textbf{does not retrain}; first 30 uploaded rows are warm-up (no prediction)
\end{itemize}
\end{column}
\begin{column}{0.45\textwidth}
\begin{exampleblock}{Row-level confidence, on hover/focus}
\small
Predicted free acid: 11.42 \\
90\% interval: 10.98–11.86 \\
Process familiarity: low \\
Confidence status: verification recommended
\end{exampleblock}
\begin{alertblock}{Caution}
\small The 90\% interval is an empirical, within-month interval. Familiarity compares inputs to the training envelope — it is \emph{not} a probability of correctness, and low familiarity does not prove the prediction is wrong.
\end{alertblock}
\end{column}
\end{columns}
\end{frame}
\note{
Two exports are available: a clean CSV with just the prediction columns, and a CSV with confidence that adds interval bounds, familiarity, confidence status, drift, and missing-feature share as separate columns — never mixed into the numeric prediction column, so it stays usable in Excel, pandas, or downstream reporting.
}

% =====================================================================
\begin{frame}{Limitations, Open Questions, and Next Steps}
\begin{columns}[T]
\begin{column}{0.5\textwidth}
\textbf{Main limitations}
\begin{itemize}
    \tiny
    \item Only one month of labeled data; adjacent minutes are dependent
    \item No independent seasonal/product-grade validation
    \item January test period has already been inspected
    \item Uncertain engineering units for several variables (acid flows, ground-phosphate, fuel, steam, wash-liquid, recycling)
    \item Process-aware ratios remain provisional; KMeans clusters are not confirmed physical regimes
    \item No causal interpretation from model importance; no control recommendation without confirmed controllable variables
    \item Empirical intervals not yet externally calibrated
    \item Target-free accuracy cannot be measured on a new CSV without some real reference values
\end{itemize}
\end{column}
\begin{column}{0.5\textwidth}
\textbf{Next steps}
\begin{enumerate}
    \tiny
    \item Obtain a later labeled month (process vars + measured free acid, ideally + grade/event info)
    \item Freeze current candidates before inspecting it
    \item Evaluate without retuning: MAE, RMSE, $R^2$, bias, transition/extreme-value errors, missing-sensor behavior, drift, interval coverage
    \item Compare the three selected candidates on that data
    \item Calibrate uncertainty intervals on genuinely later data
    \item Get process engineers to validate units, proxies, controllable variables, regimes, event windows
    \item Only then select an operational prototype
\end{enumerate}
\end{column}
\end{columns}
\end{frame}
\note{
I'd rather spend the remaining time here than on more modeling right now — the honest next step is validation data and process-engineering input, not another algorithm. I'm not proposing an LSTM, Transformer, or other larger model as the next step.
}

% =====================================================================
\begin{frame}{Conclusion}
\begin{itemize}
    \item Predicting the future \emph{change} beats persistence when the current target is available.
    \item Process-only current-value estimation (virtual sensing) is feasible within this month's data.
    \item Five-cluster Ridge is the stability-first target-free choice; the 80\% local hybrid gives the strongest retrospective balance of accuracy and stability.
    \item Just-in-time learning did not justify its added complexity and was rejected.
    \item The dashboard provides process exploration, two prediction modes, row-level reliability context, and analysis-ready exports.
    \item \textbf{The project is ready for external shadow testing, not operational control.}
    \item Next requirement: a genuinely later labeled month, plus process-engineer validation.
\end{itemize}
\end{frame}
\note{
That's where I'd like to land: solid, well-validated retrospective results on two use cases, one honest rejection, and a clear, specific ask for what I need next to move toward something operational.
}

\end{document}
